\documentclass[a4paper,11pt]{article}

\usepackage[utf8]{inputenc}

\usepackage{minted}
\usepackage{graphicx}

\begin{document}


\title{
    \textbf{Linked report}
}
\author{Jonathan Karlsson}
\date{Fall 2023}

\maketitle

\section*{Introduction}
This assignment involved creating a linked list and then adding functionality to it and benchmarking.


\section*{Implementation}
The linked list required five total methods according to the assignment. A sixth method, removeIndex was added for the stack part of the assignment.
The first method add is a method required for adding cells to the start of the sequence. It only sets the first cell as a new cell with the item as head and 
the previous first cell as the tail. 
\begin{minted}{java}
public void add(int item) {
    this.first = new Cell(item, this.first);
}
\end{minted}
The second method length works by going through every cell until the next is null and adding one to a variable named len
after every element has been counted len is returned.
\begin{minted}{java}
public int length() {
    Cell next = this.first;
    int len = 0;
    while (next != null) {
        next = next.tail;
        len++;
    }
    return len;
    }
\end{minted}
The third method, find works in a similar way to length but instead of just adding one to a total until it has gone through all elements, 
it checks each cell if its head is the same as the searched for value, if it is the method returns true, otherwise it is false.
\begin{minted}{java}
public boolean find(int item) {
    Cell next = this.first;
    while (next != null) {
        if (next.head == item)
            return true;
        next = next.tail;   
    }
    return false;
    }
\end{minted}
The fourth method, remove works in a similar way to find when finding the cell, however when it is found it instead uses a variable that saves the previous cell to remove the 
cell by setting the previous cell to point at the next cells tail. If the previous cell is null the first cell instead is pointed to next.tail. Each iteration both the previous and 
the next cell jump forwards in the list by one.
\begin{minted}{java}
public void remove(int item) {
    Cell next = this.first;
    Cell prev = null;
    while (next != null) {
        if (next.head == item) {
            if (prev != null)
                prev.tail = next.tail;
            else
                this.first = next.tail;
            break;
        }
        prev = next;
        next = next.tail;   
        
    }
}
\end{minted}
The fifth method append is very simple. It works by first using a while loop to get to the end of the current list by increasing next as long as next.tail is not
null. It then sets next.tail to the first cell of the second list and finally sets the second lists first cell to null.
\begin{minted}{java}
public void append(LinkedList b) {
    Cell next = this.first;
    while (next.tail != null) {
        next = next.tail;
    }
    next.tail = b.first;
    b.first = null;
}
\end{minted}
The sixth method removeIndex starts by moving both prev and next forward until next is at the index that should be removed. Then it checks that the index can actually be removed,
such as that next is not null which would mean that the index was one past the end of the list or that next.tail is not null, which would mean that the index is the final index, or that
prev is not null which would mean the index was 0. If the first if statement was true the value is extracted and the previous tail is set to the next tail. Else if the 
prev was null the value is extracted and the first is set to next.tail which is the second tail. The final else if says that if the next.tail is null we are at the final cell so extract the value and remove it by setting
prev.tail to null. 
\begin{minted}{java}
    public int removeIndex(int index) {
    int value = 0;
    Cell next = this.first;
    Cell prev = null;
    for (int i = 0; i<index &&next != null; i++) {
        prev = next;
        next = next.tail;
    }
    if (next != null && next.tail != null && prev != null) {
        value = next.head;
        prev.tail = next.tail;
    }
    else if (prev == null) {
        value = next.head;
        first = next.tail;
    }
    else if (next.tail == null) {
        value = next.head;
        prev.tail = null;
    }
    return value;
}
\end{minted}
The linked stack was created by using a linked list, and when pushing using the LinkedList method add to add the element
as the first cell, and when popping using the removeIndex to remove the first value and extracting its value. As the linked list can change size as it is possible 
to set the first element as the second and then create a new first to point to the second the list does not need to allocate different size as with arrays.
\begin{minted}{java}
    public void push(int value) {
        this.stack.add(value);
    }
    public int pop() {
        return this.stack.removeIndex(0);
    }    
\end{minted}
\section*{Benchmark}
Many different scenarios were benchmarked for this assignment. First was the scenario when appending a 100 element linked list to a link list of n size. 
Then appending the n sized linkedlist to the 100 element linked list. Then doing the same with arrays. The final two benchmarks were to build a linked list of n 
size and allocate an array of n size. Full table in full tables section.
\begin{table}[h]
    \begin{center}
    \scalebox{0.55}{
    \begin{tabular}{l|c|c|c|c|c|c}
    \textbf{n} & \textbf{Append 100 to n ({\textmu}s)} & \textbf{Append n to 100({\textmu}s)}& \textbf{Append 100 to n({\textmu}s)}& \textbf{Append n to 100({\textmu}s)}& \textbf{Build LinkedList({\textmu}s)}& \textbf{Allocate array({\textmu}s)}\\
    \hline
    100  &     1 &      0 &      2 &      2  &     3 &      2 \\
    200  &     1 &      0 &      2 &      3  &     4 &      2 \\
    300  &     1 &      1 &      2 &      3  &     4 &      2 \\
   1000  &     2 &      0 &      3 &      3  &     6 &      2 \\
   1500  &     2 &      0 &      3 &      3  &     8 &      2 \\
   1600  &     3 &      1 &      3 &      3  &     9 &      2 \\
  12800  &    18 &      0 &      6 &      9  &    35 &      3 \\
 102400  &   165 &      0 &     29 &     52  &   223 &     10 \\
 204800  &   317 &      0 &     60 &    103  &  7606 &     18 \\

    \end{tabular}
    }
    \caption{The result of the benchmark between linked lists and arrays}
    \label{tab:table1}
    \end{center}
\end{table}

{\noindent}It seems that the append 100 to n is slower than append n to 100, this would be because 
it is necessary to first go through the entire linked list to get to the end before appending, so when the linked list is larger
it takes longer to find the end. The arrays seem to have similar times which does make sense as both of the arrays have to be copied to a new array when appending.
This means that both the small and the large array is copied and it does not matter which is chosen first. When building the linkedlist it takes a lot longer than when allocating an array.
This is becuase the linked list has to add cells for as long as it is while the array only has to allocate a piece of memory which it will use later.
I also implemented a stack, which would definietly be faster than the previous dynamic stack which was made, this is because a stack using LinkedList can grow and shrink in size 
without having to copy any elements.
\section*{Conclusion}
Linked lists are definietly useful for some parts, for example if you want a list that is able to change its size such as the dynamic stack or when you want to append something, however they easily 
become a lot less useful whenever they grow large, as unless you only need the first element, as with a stack it will take longer and longer to get to the end of a linked list. When needing to save and work with 
large amount of data it is probably more useful to use an array, as it will be much quicker when finding the data, especially towards the end.





\newpage
\section*{Full table}
\begin{table}[h]
    \begin{center}
    \scalebox{0.55}{
    \begin{tabular}{l|c|c|c|c|c|c}
    \textbf{n} & \textbf{Append 100 to n ({\textmu}s)} & \textbf{Append n to 100({\textmu}s)}& \textbf{Append 100 to n({\textmu}s)}& \textbf{Append n to 100({\textmu}s)}& \textbf{Build LinkedList({\textmu}s)}& \textbf{Allocate array({\textmu}s)}\\
    \hline
    100  &     1 &      0 &      2 &      2  &     3 &      2 \\
    200  &     1 &      0 &      2 &      3  &     4 &      2 \\
    300  &     1 &      1 &      2 &      3  &     4 &      2 \\
    400  &     1 &      0 &      2 &      3  &     5 &      2 \\
    500  &     1 &      0 &      3 &      3  &     5 &      2 \\
    600  &     1 &      0 &      3 &      3  &     5 &      2 \\
    700  &     1 &      0 &      3 &      3  &     5 &      2 \\
    800  &     1 &      0 &      3 &      3  &     6 &      2 \\
    900  &     2 &      0 &      3 &      3  &     6 &      2 \\
   1000  &     2 &      0 &      3 &      3  &     6 &      2 \\
   1100  &     2 &      0 &      3 &      3  &     7 &      2 \\
   1200  &     2 &      0 &      3 &      3  &     7 &      2 \\
   1300  &     2 &      0 &      3 &      3  &     7 &      2 \\
   1400  &     2 &      0 &      3 &      3  &     8 &      2 \\
   1500  &     2 &      0 &      3 &      3  &     8 &      2 \\
   1600  &     3 &      1 &      3 &      3  &     9 &      2 \\
   3200  &     5 &      1 &      4 &      4  &    14 &      3 \\
   6400  &     9 &      1 &      4 &      6  &    23 &      3 \\
  12800  &    18 &      0 &      6 &      9  &    35 &      3 \\
  25600  &    37 &      0 &      9 &     15  &    64 &      4 \\
  51200  &    79 &      0 &     15 &     27  &   116 &      6 \\
 102400  &   165 &      0 &     29 &     52  &   223 &     10 \\
 204800  &   317 &      0 &     60 &    103  &  7606 &     18 \\


    \end{tabular}
    }
    \caption{The result of the benchmark between linked lists and arrays (full table)}
    \label{tab:table2}
    \end{center}
\end{table}

\end{document}
